\documentclass{amsart}
% For dealing with references we use the comment environment
\usepackage{verbatim}
\newenvironment{reference}{\comment}{\endcomment}
%\newenvironment{reference}{}{}
\newenvironment{slogan}{\comment}{\endcomment}
\newenvironment{history}{\comment}{\endcomment}

% For commutative diagrams we use Xy-pic
\usepackage[all]{xy}

% We use 2cell for 2-commutative diagrams.
\xyoption{2cell}
\UseAllTwocells

% We use multicol for the list of chapters between chapters
\usepackage{multicol}

% This is generally recommended for better output
\usepackage{lmodern}
\usepackage[T1]{fontenc}

% For cross-file-references

% Package for hypertext links:
\usepackage{hyperref}

% For any local file, say "hello.tex" you want to link to please

% Theorem environments.
%
\theoremstyle{plain}
\newtheorem{theorem}[subsection]{Theorem}
\newtheorem{proposition}[subsection]{Proposition}
\newtheorem{lemma}[subsection]{Lemma}

\theoremstyle{definition}
\newtheorem{definition}[subsection]{Definition}
\newtheorem{example}[subsection]{Example}
\newtheorem{exercise}[subsection]{Exercise}
\newtheorem{situation}[subsection]{Situation}

\theoremstyle{remark}
\newtheorem{remark}[subsection]{Remark}
\newtheorem{remarks}[subsection]{Remarks}

\numberwithin{equation}{subsection}

% Macros
%
\def\lim{\mathop{\mathrm{lim}}\nolimits}
\def\colim{\mathop{\mathrm{colim}}\nolimits}
\def\Spec{\mathop{\mathrm{Spec}}}
\def\Hom{\mathop{\mathrm{Hom}}\nolimits}
\def\Ext{\mathop{\mathrm{Ext}}\nolimits}
\def\SheafHom{\mathop{\mathcal{H}\!\mathit{om}}\nolimits}
\def\SheafExt{\mathop{\mathcal{E}\!\mathit{xt}}\nolimits}
\def\Sch{\mathit{Sch}}
\def\Mor{\mathop{\mathrm{Mor}}\nolimits}
\def\Ob{\mathop{\mathrm{Ob}}\nolimits}
\def\Sh{\mathop{\mathit{Sh}}\nolimits}
\def\NL{\mathop{N\!L}\nolimits}
\def\CH{\mathop{\mathrm{CH}}\nolimits}
\def\proetale{{pro\text{-}\acute{e}tale}}
\def\etale{{\acute{e}tale}}
\def\QCoh{\mathit{QCoh}}
\def\Ker{\mathop{\mathrm{Ker}}}
\def\Im{\mathop{\mathrm{Im}}}
\def\Coker{\mathop{\mathrm{Coker}}}
\def\Coim{\mathop{\mathrm{Coim}}}

% Boxtimes
%
\DeclareMathSymbol{\boxtimes}{\mathbin}{AMSa}{"02}

%
% Macros for moduli stacks/spaces
%
\def\QCohstack{\mathcal{QC}\!\mathit{oh}}
\def\Cohstack{\mathcal{C}\!\mathit{oh}}
\def\Spacesstack{\mathcal{S}\!\mathit{paces}}
\def\Quotfunctor{\mathrm{Quot}}
\def\Hilbfunctor{\mathrm{Hilb}}
\def\Curvesstack{\mathcal{C}\!\mathit{urves}}
\def\Polarizedstack{\mathcal{P}\!\mathit{olarized}}
\def\Complexesstack{\mathcal{C}\!\mathit{omplexes}}
% \Pic is the operator that assigns to X its picard group, usage \Pic(X)
% \Picardstack_{X/B} denotes the Picard stack of X over B
% \Picardfunctor_{X/B} denotes the Picard functor of X over B
\def\Pic{\mathop{\mathrm{Pic}}\nolimits}
\def\Picardstack{\mathcal{P}\!\mathit{ic}}
\def\Picardfunctor{\mathrm{Pic}}
\def\Deformationcategory{\mathcal{D}\!\mathit{ef}}

\setlength{\emergencystretch}{3em}
\begin{document}
\title[Stacks Project, Tag 0EAW]{445 --- Two-dimensional tame-symbol reciprocity for nonzerodivisors in Noetherian local rings}
\maketitle
\noindent Copyright (C) 2005--2025 Johan de Jong. Stacks Project authors. Source: \href{https://stacks.math.columbia.edu/tag/0EAW}{Tag 0EAW}.

\noindent Editorial difficulty note (not author text). Difficulty H3: The proof must establish cancellation over all height-one primes without assuming the local ring is regular or normal. Finite local extensions normalize the chosen factorizations, and a periodic complex with finite-length discrepancy is used to identify each tame-symbol order with a periodic Euler invariant whose sum vanishes..

\noindent Editorial provenance note (not author text). History: excerpt from revision \texttt{a0444\allowbreak 6e57e\allowbreak c1fbc\allowbreak 252a8\allowbreak 71afc\allowbreak ec775\allowbreak 2fb28\allowbreak 07b14}, chow.tex, lines 1132--1260. Reference presentation was adapted; original-author context and core support blocks were retained or added. The mathematical wording is unchanged. Licensed under GNU FDL 1.2 or later; no invariant sections or cover texts. See ../licenses/COPYING. Context blocks reproduce author definitions and supporting results; they are not additional counted problems.

\section{Tame symbols}
\label{section-tame-symbol}

\noindent
Consider a Noetherian local ring $(A, \mathfrak m)$ of dimension $1$.
We denote $Q(A)$ the total ring of fractions of $A$, see
Algebra, Example \href{https://stacks.math.columbia.edu/tag/02C5}{example localize at prime (Tag 02C5)}.
The {\it tame symbol} will be a map
$$
\partial_A(-, -) : Q(A)^* \times Q(A)^* \longrightarrow \kappa(\mathfrak m)^*
$$
satisfying the following properties:
\begin{enumerate}
\item $\partial_A(f, gh) = \partial_A(f, g) \partial_A(f, h)$
\label{item-bilinear}
for $f, g, h \in Q(A)^*$,
\item $\partial_A(f, g) \partial_A(g, f) = 1$
\label{item-skew}
for $f, g \in Q(A)^*$,
\item $\partial_A(f, 1 - f) = 1$
\label{item-1-x}
for $f \in Q(A)^*$ such that $1 - f \in Q(A)^*$,
\item $\partial_A(aa', b) = \partial_A(a, b)\partial_A(a', b)$
\label{item-bilinear-better}
and $\partial_A(a, bb') = \partial_A(a, b)\partial_A(a, b')$
for $a, a', b, b' \in A$ nonzerodivisors,
\item $\partial_A(b, b) = (-1)^m$
\label{item-skew-better}
with $m = \text{length}_A(A/bA)$
for $b \in A$ a nonzerodivisor,
\item $\partial_A(u, b) = u^m \bmod \mathfrak m$
\label{item-normalization}
with $m = \text{length}_A(A/bA)$ for $u \in A$ a unit and
$b \in A$ a nonzerodivisor, and
\item
\label{item-1-x-better}
$\partial_A(a, b - a)\partial_A(b, b) = \partial_A(b, b - a)\partial_A(a, b)$
for $a, b \in A$ such that $a, b, b - a$ are nonzerodivisors.
\end{enumerate}
Since it is easier to work with elements of $A$ we will
often think of $\partial_A$ as a map defined on pairs of
nonzerodivisors of $A$ satisfying (\ref{item-bilinear-better}),
(\ref{item-skew-better}), (\ref{item-normalization}),
(\ref{item-1-x-better}). It is an exercise to see that
setting
$$
\partial_A(\frac{a}{b}, \frac{c}{d}) =
\partial_A(a, c) \partial_A(a, d)^{-1} \partial_A(b, c)^{-1} \partial_A(b, d)
$$
we get a well defined map $Q(A)^* \times Q(A)^* \to \kappa(\mathfrak m)^*$
satisfying (\ref{item-bilinear}), (\ref{item-skew}), (\ref{item-1-x})
as well as the other properties.

\medskip\noindent
We do not claim there is a unique map with these properties.
Instead, we will give a recipe for constructing such a map.
Namely, given $a_1, a_2 \in A$ nonzerodivisors, we choose
a ring extension $A \subset B$ and local factorizations
as in Lemma \href{https://stacks.math.columbia.edu/tag/0EAG}{lemma prepare tame symbol (Tag 0EAG)}.
Then we define
\begin{equation}
\label{equation-tame-symbol}
\partial_A(a_1, a_2) = \prod\nolimits_j
\text{Norm}_{\kappa(\mathfrak m_j)/\kappa(\mathfrak m)}
((-1)^{e_{1, j}e_{2, j}}u_{1, j}^{e_{2, j}}u_{2, j}^{-e_{1, j}}
\bmod \mathfrak m_j)^{m_j}
\end{equation}
where $m_j = \text{length}_{B_j}(B_j/\pi_j B_j)$ and the product
is taken over the maximal ideals $\mathfrak m_1, \ldots, \mathfrak m_r$ of $B$.



\begin{lemma}
\label{lemma-key-nonzerodivisors}
Let $(A, \mathfrak m)$ be a $2$-dimensional Noetherian local ring.
Let $a, b \in A$ be nonzerodivisors.
Then we have
$$
\sum
\text{ord}_{A/\mathfrak q}(\partial_{A_{\mathfrak q}}(a, b))
=
0
$$
where the sum is over the height $1$ primes $\mathfrak q$ of $A$.
\end{lemma}

\begin{proof}
If $\mathfrak q$ is a height $1$ prime of $A$ such that $a, b$
map to a unit of $A_\mathfrak q$, then $\partial_{A_\mathfrak q}(a, b) = 1$.
Thus the sum is finite. In fact, if
$V(ab) = \{\mathfrak m, \mathfrak q_1, \ldots, \mathfrak q_r\}$
then the sum is over $i = 1, \ldots, r$.
For each $i$ we pick an extension $A_{\mathfrak q_i} \subset B_i$
as in Lemma \href{https://stacks.math.columbia.edu/tag/0EAG}{lemma prepare tame symbol (Tag 0EAG)} for $a, b$.
By Lemma \href{https://stacks.math.columbia.edu/tag/0EAV}{lemma perpare key (Tag 0EAV)} with $t = ab$ and the given list of primes
we may assume we have a finite local extension $A \subset B$
with $B/A$ annihilated by a power of $ab$ and such that
for each $i$ the $B_{\mathfrak q_i} \cong B_i$.
Observe that if $\mathfrak q_{i, j}$ are the primes of $B$
lying over $\mathfrak q_i$ then we have
$$
\text{ord}_{A/\mathfrak q_i}(\partial_{A_{\mathfrak q_i}}(a, b))
=
\sum\nolimits_j
\text{ord}_{B/\mathfrak q_{i, j}}(\partial_{B_{\mathfrak q_{i, j}}}(a, b))
$$
by Lemma \href{https://stacks.math.columbia.edu/tag/0EAT}{lemma norm down tame symbol (Tag 0EAT)} and
Algebra, Lemma \href{https://stacks.math.columbia.edu/tag/02MJ}{lemma finite extension dim 1 (Tag 02MJ)}.
Thus we may replace $A$ by $B$ and
reduce to the case discussed in the next paragraph.

\medskip\noindent
Assume for each $i$ there is a nonzerodivisor
$\pi_i \in A_{\mathfrak q_i}$ and units $u_i, v_i \in A_{\mathfrak q_i}$
such that for some integers $e_i, f_i \geq 0$ we have
$$
a = u_i \pi_i^{e_i},\quad b = v_i \pi_i^{f_i}
$$
in $A_{\mathfrak q_i}$. Setting
$m_i = \text{length}_{A_{\mathfrak q_i}}(A_{\mathfrak q_i}/\pi_i)$
we have
$\partial_{A_{\mathfrak q_i}}(a, b) =
((-1)^{e_if_i}u_i^{f_i}v_i^{-e_i})^{m_i}$ by definition.
Since $a, b$ are nonzerodivisors the
$(2, 1)$-periodic complex $(A/(ab), a, b)$ has vanishing cohomology.
Denote $M_i$ the image of $A/(ab)$ in $A_{\mathfrak q_i}/(ab)$.
Then we have a map
$$
A/(ab) \longrightarrow \bigoplus M_i
$$
whose kernel and cokernel are supported in $\{\mathfrak m\}$
and hence have finite length. Thus we see that
$$
\sum e_A(M_i, a, b) = 0
$$
by Lemma \href{https://stacks.math.columbia.edu/tag/0EA9}{lemma compare periodic lengths (Tag 0EA9)}. Hence it suffices to show
$e_A(M_i, a, b) =
- \text{ord}_{A/\mathfrak q_i}(\partial_{A_{\mathfrak q_i}}(a, b))$.

\medskip\noindent
Let us prove this first, in case $\pi_i, u_i, v_i$ are the images
of elements $\pi_i, u_i, v_i \in A$ (using the same symbols
should not cause any confusion). In this case we get
\begin{align*}
e_A(M_i, a, b) & =
e_A(M_i, u_i\pi_i^{e_i}, v_i\pi_i^{f_i}) \\
& =
e_A(M_i, \pi_i^{e_i}, \pi_i^{f_i}) -
e_A(\pi_i^{e_i}M_i, 0, u_i) +
e_A(\pi_i^{f_i}M_i, 0, v_i) \\
& =
0 -
f_im_i\text{ord}_{A/\mathfrak q_i}(u_i) +
e_im_i\text{ord}_{A/\mathfrak q_i}(v_i) \\
& =
-m_i\text{ord}_{A/\mathfrak q_i}(u_i^{f_i}v_i^{-e_i}) =
-\text{ord}_{A/\mathfrak q_i}(\partial_{A_{\mathfrak q_i}}(a, b))
\end{align*}
The second equality holds by Lemma \href{https://stacks.math.columbia.edu/tag/0EAC}{lemma multiply period length (Tag 0EAC)}.
Observe that
$M_i \subset (M_i)_{\mathfrak q_i} = A_{\mathfrak q_i}/(\pi_i^{e_i + f_i})$
and
$(\pi_i^{e_i}M_i)_{\mathfrak q_i} \cong A_{\mathfrak q_i}/\pi_i^{f_i}$ and
$(\pi_i^{f_i}M_i)_{\mathfrak q_i} \cong A_{\mathfrak q_i}/\pi_i^{e_i}$.
The $0$ in the third equality comes from
Lemma \href{https://stacks.math.columbia.edu/tag/0EAB}{lemma powers period length zero (Tag 0EAB)}
and the other two terms come from
Lemma \href{https://stacks.math.columbia.edu/tag/02QF}{lemma length multiplication (Tag 02QF)}.
The last two equalities follow from multiplicativity of
the order function and from the definition of our tame symbol.

\medskip\noindent
In general, we may first choose $c \in A$, $c \not \in \mathfrak q_i$
such that $c\pi_i \in A$. After replacing $\pi_i$ by $c\pi_i$
and $u_i$ by $c^{-e_i}u_i$ and $v_i$ by $c^{-f_i}v_i$
we may and do assume $\pi_i$ is in $A$.
Next, choose an $c \in A$, $c \not \in \mathfrak q_i$
with $cu_i, cv_i \in A$. Then we observe that
$$
e_A(M_i, ca, cb) = e_A(M_i, a, b) - e_A(aM_i, 0, c) + e_A(bM_i, 0, c)
$$
by Lemma \href{https://stacks.math.columbia.edu/tag/02QF}{lemma length multiplication (Tag 02QF)}.
On the other hand, we have
$$
\partial_{A_{\mathfrak q_i}}(ca, cb) =
c^{m_i(f_i - e_i)}\partial_{A_{\mathfrak q_i}}(a, b)
$$
in $\kappa(\mathfrak q_i)^*$ because $c$ is a unit in $A_{\mathfrak q_i}$.
The arguments in the previous paragraph show that
$e_A(M_i, ca, cb) = -
\text{ord}_{A/\mathfrak q_i}(\partial_{A_{\mathfrak q_i}}(ca, cb))$.
Thus it suffices to prove
$$
e_A(aM_i, 0, c) = \text{ord}_{A/\mathfrak q_i}(c^{m_if_i})
\quad\text{and}\quad
e_A(bM_i, 0, c) = \text{ord}_{A/\mathfrak q_i}(c^{m_ie_i})
$$
and this follows from Lemma \href{https://stacks.math.columbia.edu/tag/02QF}{lemma length multiplication (Tag 02QF)}
by the description (see above)
of what happens when we localize at $\mathfrak q_i$.
\end{proof}
\end{document}
